% ===================================================================== Chapter 13
\chapter{Buckling Analysis}\label{ch:buckling}

\section{Why Buckling Was Investigated}
The LC2 static state stores a compressive force of 548.94~kN in a column with a slenderness ratio $L/r$ of 53.67. A static stress check
alone cannot tell whether that state is stable: a yield utilisation of 0.578 says nothing about the load at which the straight equilibrium
bifurcates into a laterally deflected one. Because the LC2 restraint is exactly the condition that creates the compressive force, its
stability had to be examined before the static result could be interpreted.

\section{Applied Thermal Compression}
The column load is the LC2 end reaction, $N$ = 548,936.6~N, equivalent to a mean axial stress of 582.44~MPa over the section area of
942.48~mm$^2$. The 605.16~MPa von Mises peak is a local stress, not the column load. The load is thermal and displacement-controlled: it is
the force needed to hold the thermally expanded duct at its original length. In the linear buckling analysis the load factor $\lambda$
multiplies this entire thermal pre-stress; $\lambda$ = 1.108 therefore corresponds to about 10.8\,\% more restrained thermal strain, or
roughly 24~K more mean temperature rise.

\section{Euler Estimate}
Hand estimates were made before the FE analysis (Table~\ref{tab:section}), using Equation~\eqref{eq:euler} with a Rayleigh average of the
axially varying modulus $E(z)$ and, separately, with the shear correction. For the guided end condition of the LC2 supports ($K$ = 1) they give
load factors of 1.119 (Euler) and 1.104 (Euler + shear). All end conditions give intermediate columns ($KL/r$ below $C_c$ = 60.2), and the
elastic Euler stress of about 650~MPa exceeds the conventional proportional limit. The conventional Johnson estimate for the guided column
is 1.058, and the squash load corresponds to a factor of 1.755. Local or shell buckling is not credible: with $R/t$ = 1.5 the tube is far
outside thin-shell theory and compact by the round-HSS criterion $D/t < 0.11E/F_y$ ($4 < 20.2$) \cite{aisc360}.

%% generated by gen_tables.py - RE-ANALYSIS 2026
\begin{table}[htbp]
\centering
\footnotesize
\caption{Column properties and hand estimates of the buckling load factor (Section 8A)}
\label{tab:section}
\begin{tabular}{lrrr}
\toprule
End condition & $KL/r$ & Euler, $E(z)$ & Euler + shear \\
\midrule
Guided (rotation fixed, sway free) = S1 & 53.7 & 1.119 & 1.104 \\
Pinned--pinned & 53.7 & 1.114 & 1.099 \\
Fixed--pinned & 37.5 & 2.29 & 2.22 \\
Fixed--fixed (no sway) & 26.8 & 4.47 & 4.23 \\
\midrule
\multicolumn{4}{l}{$A$ = 942.48 mm$^2$; $I$ = 117,810 mm$^4$; $r$ = 11.18 mm; $L/r$ = 53.67; $C_c$ = 60.2; Johnson (guided) 1.058; squash 1.755} \\
\bottomrule
\end{tabular}
\par\smallskip\parbox{\linewidth}{\footnotesize\raggedright Source: \texttt{08\_Structural\_Analysis/Buckling/BUCKLING\_RESULTS.md} and \texttt{PROJECT\_STATE.md} \S15.1. Load factor = $P_{cr}/N$ with $N$ = 548,936.6 N. The shear correction uses the Engesser form with Cowper's coefficient $k$ = 0.620 \cite{cowper1966}.}
\end{table}


\section{ANSYS Linear Buckling}
The Eigenvalue Buckling system was linked to the LC2 static solution in a save-as copy of the project, as a perturbation restart that keeps
all LC2 boundary conditions. The re-solved LC2 pre-stress in the copy reproduced the Section 7B values exactly (peak von Mises and end
reaction). The full 108,252-node model solved within the Student licence. Before the project model was solved, the workflow was benchmarked
on a constant-property toy tube whose Euler loads are known in closed form for both end conditions.

\section{Mode Shapes}
The first mode (Figure~\ref{fig:bkmode}) is a global column mode: the end planes stay perpendicular to the axis but translate sideways in
opposite directions; the lateral deflection is zero at mid-span and the bending curvature is largest at the rotation-held end sections. The
cross-section moves as a rigid section (99.9999\,\%); there is no ovalisation, shell or local mode. The shape correlates with
$\cos(\pi z/L)$ to 0.999997. Modes 1 and 2 are an orthogonal pair with the same eigenvalue, as expected for an
axisymmetric column.

\begin{figure}[htbp]
\centering
\begin{subfigure}[t]{0.49\linewidth}\centering\includegraphics[width=\linewidth,height=52mm,keepaspectratio]{LC2_Linear_Buckling_Mode_1_Total_Deformation_iso.png}\caption{Isometric view}\end{subfigure}\hfill
\begin{subfigure}[t]{0.49\linewidth}\centering\includegraphics[width=\linewidth,height=52mm,keepaspectratio]{LC2_Linear_Buckling_Mode_1_Total_Deformation_side_YZ.png}\caption{Side view}\end{subfigure}
\caption[First buckling mode of the restrained duct under S1]{First buckling mode of the restrained duct under S1 (load multiplier 1.108; normalised eigenvector, amplitude arbitrary).}
\label{fig:bkmode}
\src{\provmech}
\end{figure}


\section{Buckling Factor}
The first eigenvalue is $\lambda_1$ = 1.108 (Table~\ref{tab:buckling}), giving a critical load $P_{cr} = \lambda_1 N$ = 608.25~kN. The next
eigenvalue pair is 4.298. The FE value lies between the Euler estimate with $E(z)$ ($-$1.0\,\%) and the shear-corrected estimate (+0.4\,\%),
which is the expected position for a stocky tube (Figure~\ref{fig:bkpcr}).

%% generated by gen_tables.py - RE-ANALYSIS 2026
\begin{table}[htbp]
\centering
\footnotesize
\caption{Linear eigenvalue buckling results of the baseline under S1 (Section 8A)}
\label{tab:buckling}
\begin{tabular}{>{\raggedright\arraybackslash}p{0.45\linewidth}r}
\toprule
Quantity & Value \\
\midrule
Applied axial compression $N$ (LC2 end reaction) [N] & 548,936.6 \\
$\lambda_1 = \lambda_2$ (orthogonal pair) & 1.10805 \\
Next eigenvalue pair & 4.298 \\
Critical load $P_{cr} = \lambda_1 N$ [kN] & 608.25 \\
Dominant mode & global guided-column sway, $\cos(\pi z/L)$ \\
First-yield load factor (LC2) & 1.730 \\
Conventional inelastic (Johnson) estimate & 1.058 \\
\bottomrule
\end{tabular}
\par\smallskip\parbox{\linewidth}{\footnotesize\raggedright Source: \texttt{MASTER\_PROJECT\_DATA.csv} rows M091, M096, M098--M100, M107; \texttt{BUCKLING\_RESULTS.md}. $\lambda_1$ is an idealised linear stability indicator, not a collapse load and not a factor of safety.}
\end{table}


\begin{figure}[htbp]
\centering
\begin{subfigure}[t]{0.49\linewidth}\centering\includegraphics[width=\linewidth,height=52mm,keepaspectratio]{F8A_03_critical_load_summary.png}\caption{Critical axial force [kN]}\end{subfigure}\hfill
\begin{subfigure}[t]{0.49\linewidth}\centering\includegraphics[width=\linewidth,height=52mm,keepaspectratio]{F8A_04_load_factor_comparison.png}\caption{Buckling load factors}\end{subfigure}
\caption[Hand estimates and FE results against the applied restrained force]{Hand estimates and FE results against the applied restrained force: elastic Euler, shear-corrected, conventional Johnson (inelastic) and FE linear buckling.}
\label{fig:bkpcr}
\src{\provdata}
\end{figure}

\section{Buckling Mesh Sensitivity}
The same buckling analysis was repeated on the six meshes of Chapter~\ref{ch:smesh} that carry a buckling solution (Table~\ref{tab:appF}). $\lambda_1$ ranges from 1.1079876 (XC) to 1.1080483 (FA); the coarsening family is monotonic with a GCI of
$5.6\times10^{-6}$, and the fine meshes lie within $6\times10^{-6}$ of mesh B. The mode is the same guided-sway shape on every mesh. The buckling
result is not sensitive to the structural mesh.


\section{Yield vs Buckling}
Two different load factors describe the idealised LC2 state and are reported side by side:
\begin{itemize}
\item the first-yield factor, 1.730: the thermal load multiplier at which the most highly stressed point reaches the local yield strength;
\item the linear buckling factor, $\lambda_1$ = 1.108: the multiplier at which the perfect, elastic column bifurcates.
\end{itemize}
Because $\lambda_1$ is below the first-yield factor, the idealised S1 model is stability-limited: elastic bifurcation is reached before first
yield. With $\lambda_1$ this close to 1, the idealised model is close to instability at the applied thermal load. The yield utilisation of
0.578 is therefore not a structural margin. Expressed on the Section 2 temperature basis (254.7~K mean rise instead of 225.5~K), the same
analysis would give $\lambda_1 \approx$ 0.98. No single factor of safety is formed in this project.

\section{Limitations of Linear Eigenvalue Buckling}
$\lambda_1$ is a linear stability indicator for an idealised elastic model. It is \textbf{not} a guaranteed real-world collapse load. The
analysis assumes:
\begin{itemize}
\item a perfectly straight tube without out-of-straightness, load eccentricity or residual stress;
\item linear elastic material, although the columns are intermediate and the elastic buckling stress exceeds the proportional limit;
\item small deflections, with no post-buckling path;
\item idealised supports whose relation to a real flange is undefined.
\end{itemize}
A real, imperfect duct would deflect laterally at loads below $\lambda_1$, and inelastic behaviour lowers its capacity further; the
conventional Johnson estimate of 1.058 indicates the direction of that reduction, not its material-specific size. Rotational flexibility of
real end fixtures would lower $\lambda_1$ below the S1 value, so S1 is not a lower bound. An imperfection-sensitive, geometrically and
materially nonlinear analysis would be required to determine a collapse load (Chapter~\ref{ch:future}). An elastic, geometrically
nonlinear imperfection-sensitivity extension is reported in Section~\ref{sec:lsdyna}; it does not determine a collapse load.
